\subsection{False-claim soundness and repeated-opening failures}\label{sec:zkrs:soundness}

\begin{lemma}[False claim forces a partial-evaluation discrepancy]\label{lem:zkrs:false-error}
Fix a Boolean table $w$, packed columns $\widehat w$,
and evaluation points $(r_{\rm lo},r_{\rm hi})$.
Suppose the public $y\neq y_w$.
If the claim and incoming affine equations pass,
then
\[
 e:=m+s^*+p\in L^I
\]
is nonzero.
\end{lemma}
\begin{proof}
If $e=0$ then $m=s^*+p$.
The claim and incoming equations imply
$y=\sum_vq_v(m_v+p_v)$ by
Lemma~\ref{lem:pack:incoming-eliminate}.
Hence $y=\sum_vq_vs_v^*=y_w$ by
Lemma~\ref{lem:pack:partial-eval},
contradicting the false claim.
\end{proof}

\begin{lemma}[Correct hidden transposition gives a nonzero outgoing polynomial]\label{lem:zkrs:correct-c}
Under Lemma~\ref{lem:zkrs:false-error},
if $c=c^p$, then passing the outgoing
equation is equivalent to $B_e(\theta)=0$,
where $e$ is fixed and nonzero.
\end{lemma}
\begin{proof}
By Lemma~\ref{lem:pack:out-error},
the outgoing discrepancy is
$B_e(\theta)+D_{c+c^p}(\theta)$.
The assumed equality $c=c^p$
makes the second term zero.
By Lemma~\ref{lem:zkrs:false-error},
$e\neq0$, and the outgoing
equation holds precisely when
the first term vanishes.
\end{proof}

\begin{lemma}[Incorrect hidden transposition gives a degree-$\ell$ mismatch]\label{lem:zkrs:incorrect-c}
If $c\neq c^p$ is fixed before a uniform
$\theta\in L^\ell$, the probability that
$D_c(\theta)=\Delta_p(\theta)$ is at most
$\min\{1,\ell/q\}$.
\end{lemma}
\begin{proof}
Apply
Corollary~\ref{cor:binary:mismatch-bound}.
Its hypotheses are exactly the stated
coefficient mismatch and independence
from the challenge.
\end{proof}

\begin{theorem}[Ideal masked ring-switching false-claim soundness]\label{thm:zkrs:false-sound}
Fix a Boolean witness, its exact packed
columns, public $(r_{\rm lo},r_{\rm hi},y)$,
and the public $m,y'$ and hidden
$p,p_{\rm in},c$ before the independent
uniform first challenge $\theta$.
After $\theta$, a malicious prover may
choose the checker values $(t,\rho,\sigma,U)$
adaptively, but must bind them before the
independent uniform second challenge $z$.
It may choose $A$ after receiving $z$.
Assume all ideal field-affine checks and
public trace checks of
Definition~\ref{def:zkrs:affine-checks}
are enforced exactly.
If $y\ne y_w$, then
\[
 \boxed{
 \Pr[\mathrm{accept}]
 \leq\min\left\{1,\frac{\ell+n}{q}\right\}.
 }
\]
\end{theorem}
\begin{proof}
Distinguish two mutually exclusive cases,
which are already fixed before $\theta$.

If $c=c^p$, whenever the claim and
incoming equations pass the vector
$e=m+s^*+p$ is nonzero by
Lemma~\ref{lem:zkrs:false-error}.
By Lemma~\ref{lem:zkrs:correct-c},
outgoing acceptance requires
$B_e(\theta)=0$. The discrepancy
polynomial has degree at most $\ell$
and is nonzero; Theorem~\ref{thm:pack:batch-sound}
bounds this event by $\ell/q$.
Further checker constraints can only
reduce acceptance.

If $c\ne c^p$, let $E$ denote the
exceptional first-challenge event
$D_c(\theta)=\Delta_p(\theta)$.
Lemma~\ref{lem:zkrs:incorrect-c}
gives $\Pr[E]\leq\ell/q$.
On its complement, if the public
trace checks pass then
Lemma~\ref{lem:binary:forced-residual}
guarantees a nonzero vector
$e_k=U_k+V_k+t_k^2+t_k$
for at least one coordinate.
The auxiliary values are fixed
before the fresh $z$, so the residual
detection theorem
(Theorem~\ref{thm:zkrs:residual-detection})
bounds passage of the two anchored
equations by $n/q$.
Conditioning on the malicious prover's
chosen state and private random tape
and then averaging gives
$\Pr[\mathrm{accept}]\leq(\ell+n)/q$
in this case.
Taking the maximum of the two bounds,
and capping by one, proves the theorem.
\end{proof}

\begin{definition}[Two-opening pad-reuse experiment]\label{def:zkrs:reuse-game}
Let a single witness $w$ be evaluated
at two distinct high-axis points
$r_{\rm hi}^{(1)},r_{\rm hi}^{(2)}$
with the same low-axis point.
Let $s^{*(i)}$ be the associated
partial-evaluation vectors.
A \emph{reused-pad} implementation
chooses one uniform $p\in L^I$
and publishes
$m^{(1)}=s^{*(1)}+p$ and
$m^{(2)}=s^{*(2)}+p$.
\end{definition}

\begin{lemma}[Reusing a pad reveals the difference of partial evaluations]\label{lem:zkrs:reuse-difference}
In the two-opening pad-reuse experiment,
\[
 \boxed{m^{(1)}+m^{(2)}
        =s^{*(1)}+s^{*(2)}}
\]
for every outcome of the pad.
\end{lemma}
\begin{proof}
Add the two definitions:
$m^{(1)}+m^{(2)}
=s^{*(1)}+p+s^{*(2)}+p$.
Because $p+p=0$ in characteristic
two, the pad cancels.
\end{proof}

\begin{theorem}[Perfect two-opening leakage under pad reuse]\label{thm:zkrs:reuse-attack}
For any $\ell\geq1$ and $h\geq1$,
there are two Boolean witnesses having the
same authorized evaluation values at two
public points, but whose reused-pad
public pairs $(m^{(1)},m^{(2)})$
have statistical distance $1$.
\end{theorem}
\begin{proof}
Take $r_{\rm lo}=0_\ell$ and
$r_{\rm hi}^{(1)}=0_h$,
$r_{\rm hi}^{(2)}=(1,0,\ldots,0)$.
Let $w^{(0)}$ be identically zero.
Choose $v_*\neq0_\ell$ and define
$w^{(1)}$ with exactly one nonzero
entry at $(v_*,0_h)$.
Both public evaluations, at both
chosen points, are zero for both
witnesses: the low-axis equality
weights select only the row
$v=0_\ell$, which is zero.

For $w^{(0)}$, both partial vectors
vanish, so Lemma~\ref{lem:zkrs:reuse-difference}
gives $m^{(1)}+m^{(2)}=0$ always.
For $w^{(1)}$, the first high-axis
point selects the Boolean column
$j=0_h$, producing a partial vector
with a single $1$ in row $v_*$.
The second high-axis point selects
the column $(1,0,\ldots,0)$, where
the witness is zero, and hence
$s^{*(2)}=0$.
Thus $m^{(1)}+m^{(2)}$ has a nonzero
entry at $v_*$ with probability one.
The deterministic event
$m^{(1)}+m^{(2)}=0$ distinguishes the
two distributions with probabilities
$1$ and $0$.
Lemma~\ref{lem:security:events}
makes the statistical distance $1$.
\end{proof}

\begin{corollary}[Fresh partial-evaluation pads are necessary for this implementation]\label{cor:zkrs:fresh-pads}
Any multi-opening version that
publishes raw masked partial
vectors in the form $m=s^*+p$
must not reuse the same pad
across arbitrary evaluation points
if it is to satisfy witness privacy
for the authorized evaluation values.
\end{corollary}
\begin{proof}
The reuse construction of
Definition~\ref{def:zkrs:reuse-game}
violates the two-witness closeness
required by statistical zero
knowledge, as shown in
Theorem~\ref{thm:zkrs:reuse-attack}.
Thus that reuse rule is not
compatible with the stated
privacy goal.
\end{proof}

\begin{remark}[No universal lower bound]
Corollary~\ref{cor:zkrs:fresh-pads}
rules out this exact pad-reuse pattern;
it does not prove an information-theoretic
lower bound for all possible multi-opening
ZK schemes. Correlated masks with
additional zero-knowledge machinery
would require their own proof.
\end{remark}
